\documentclass[../DoAn.tex]{subfiles}
\begin{document}

Phụ lục này trình bày đặc tả chi tiết bốn Use Case cốt lõi của phân hệ Quản lý Thông tin Nhiệm vụ ROV.

\section{UC-01: Tải lên và đồng bộ dữ liệu cảm biến}

\begin{longtable}{|p{3.5cm}|p{10.0cm}|}
\caption{Đặc tả Use Case UC-01: Tải lên và đồng bộ dữ liệu cảm biến}\label{tab:b-uc01}\\
\hline
\textbf{Mã Use Case} & UC-01 \\
\hline
\textbf{Tên} & Tải lên và đồng bộ dữ liệu cảm biến (Sensor Data Sync) \\
\hline
\textbf{Tác nhân} & Operator \\
\hline
\textbf{Tiền điều kiện} & Tác nhân đã đăng nhập và đang ở trang chi tiết Lượt lặn (TripDetailPage). \\
\hline
\textbf{Luồng chính} &
(1) Tác nhân kéo thả hoặc chọn tệp CSV/DVL/Sonar vào khu vực tải lên. \newline
(2) Hệ thống xác thực định dạng tệp (CSV phải chứa hai cột bắt buộc: \texttt{timestamp} và \texttt{depth}; các cột còn lại như \texttt{temp}, \texttt{pressure}, \texttt{yaw}, \texttt{voltage},... là tùy chọn). \newline
(3) Hệ thống bóc tách tọa độ GPS từ dòng đầu tiên, gọi API Nominatim để chuyển thành tên địa danh (\textit{Reverse Geocoding}). \newline
(4) Hệ thống lưu dữ liệu vào MongoDB, tính toán phát hiện bất thường bằng Z-Score ($|z| > 2{,}5$). \newline
(5) Giao diện hiển thị biểu đồ môi trường; khi video phát, thanh đánh dấu dọc di chuyển đồng bộ theo dấu thời gian. \\
\hline
\textbf{Luồng phát sinh} & Định dạng tệp không hợp lệ $\to$ hệ thống thông báo lỗi chi tiết và yêu cầu tải lại. Tệp trùng phạm vi thời gian $\to$ hệ thống cảnh báo số dòng bị bỏ qua (\texttt{overlap\_trimmed}). \\
\hline
\textbf{Hậu điều kiện} & Dữ liệu cảm biến được liên kết thành công vào lượt lặn; biểu đồ và bản đồ GPS cập nhật tức thì. \\
\hline
\end{longtable}

\section{UC-02: Kích hoạt phân tích nhận diện vật thể bằng YOLOv8}

\begin{longtable}{|p{3.5cm}|p{10.0cm}|}
\caption{Đặc tả Use Case UC-02: Phân tích nhận diện vật thể bằng AI}\label{tab:b-uc02}\\
\hline
\textbf{Mã Use Case} & UC-02 \\
\hline
\textbf{Tên} & Kích hoạt phân tích nhận diện vật thể (AI Object Detection) \\
\hline
\textbf{Tác nhân} & Operator; YOLO Service (hệ thống bên ngoài) \\
\hline
\textbf{Tiền điều kiện} & Tệp video hoặc ảnh đã được tải lên thành công (trạng thái \textit{Ready}). \\
\hline
\textbf{Luồng chính} &
(1) Tác nhân nhấn nút Cài đặt AI trên thanh công cụ video; cửa sổ thiết lập hiện ra cho phép chọn mô hình (YOLOv8n/tùy chỉnh) và ngưỡng độ tin cậy. \newline
(2) Tác nhân nhấn \textit{``Chạy phân tích''}; hệ thống đưa tác vụ vào hàng đợi Bull Queue và trả về xác nhận tiếp nhận (HTTP 202). \newline
(3) Worker nền gọi Python Microservice, truyền URL presigned của tệp phương tiện, mô hình và ngưỡng. \newline
(4) YOLO Service xử lý từng khung hình (video: lấy mẫu theo tần suất thích nghi), trả về danh sách nhãn kèm tọa độ bounding box và thời điểm khung hình (\texttt{frameTime}). \newline
(5) Hệ thống lưu kết quả vào MongoDB, đẩy thông báo SSE đến giao diện tác nhân. \newline
(6) Bounding box xuất hiện trên video đồng bộ theo từng khung hình đang phát. \\
\hline
\textbf{Luồng phát sinh} & YOLO Service không phản hồi hoặc quá tải $\to$ job thất bại sau 1 lần thử (timeout 25 phút); Worker ghi \texttt{analysisStatus = 'failed'} và thông báo lỗi qua SSE. \\
\hline
\textbf{Hậu điều kiện} & Nhãn nhận diện (Labels, Bbox, FrameTime, TrackId) được lưu và hiển thị theo thời gian thực trên giao diện. \\
\hline
\end{longtable}

\section{UC-03: Trích xuất và phân tích bằng chứng}

\begin{longtable}{|p{3.5cm}|p{10.0cm}|}
\caption{Đặc tả Use Case UC-03: Trích xuất bằng chứng (Evidence System)}\label{tab:b-uc03}\\
\hline
\textbf{Mã Use Case} & UC-03 \\
\hline
\textbf{Tên} & Trích xuất và phân tích bằng chứng (Evidence System) \\
\hline
\textbf{Tác nhân} & Operator \\
\hline
\textbf{Tiền điều kiện} & Video của lượt lặn đang được phát trên TripDetailPage. \\
\hline
\textbf{Luồng chính} &
(1) Tác nhân phát hiện vật thể nghi vấn, nhấn \textit{Photo} hoặc \textit{Clip} trên thanh công cụ. \newline
(2a) \textit{Photo:} Hệ thống trích xuất khung hình hiện tại thành tệp PNG, tải lên AWS S3. Nếu bounding box đang hiển thị, tọa độ được ghi vào ảnh (canvas burn-in). \newline
(2b) \textit{Clip:} Tác nhân nhấn bắt đầu $\to$ dừng; hệ thống ghi nhận $t_{\text{start}}$ và $t_{\text{end}}$ mà không cắt tệp vật lý. \newline
(3) Bằng chứng được lưu vào collection \textit{Snapshots}; tác nhân có thể ghi chú thêm. \newline
(4) Tác nhân tùy chọn kích hoạt YOLO riêng cho bằng chứng này (mở rộng từ UC-02). \\
\hline
\textbf{Luồng phát sinh} & Không có video đang active $\to$ nút bị vô hiệu hóa, không thể thao tác. \\
\hline
\textbf{Hậu điều kiện} & Bằng chứng được lưu thành công và hiển thị trong bảng Evidence; tệp video gốc không bị thay đổi. \\
\hline
\end{longtable}

\section{UC-04: Tạo báo cáo tóm tắt bằng ngôn ngữ tự nhiên}

\begin{longtable}{|p{3.5cm}|p{10.0cm}|}
\caption{Đặc tả Use Case UC-04: Tóm tắt nhiệm vụ bằng AI (AI Summary)}\label{tab:b-uc04}\\
\hline
\textbf{Mã Use Case} & UC-04 \\
\hline
\textbf{Tên} & Tạo báo cáo tóm tắt nhiệm vụ bằng AI (AI Summary) \\
\hline
\textbf{Tác nhân} & Operator; Gemini AI (hệ thống bên ngoài) \\
\hline
\textbf{Tiền điều kiện} & Trạng thái Chuyến khảo sát (Project) đã chuyển sang \textit{Completed}. \\
\hline
\textbf{Luồng chính} &
(1) Tác nhân truy cập trang ProjectDetailPage, nhấn nút \textit{``Tạo Tóm tắt''}. \newline
(2) Hệ thống thu thập thống kê chuyến đi: địa điểm (\texttt{locationName}), thời gian, số lượt lặn, tổng số ảnh/video, số cảnh báo cảm biến (Z-Score $|z| > 2{,}5$). \newline
(3) Dữ liệu được đưa vào mẫu câu (Prompt) và gửi tới Gemini API qua Bull Queue. \newline
(4) Hệ thống nhận phản hồi, lưu văn bản tóm tắt vào trường \texttt{project.aiSummary}. \newline
(5) Thông báo SSE được đẩy đến người kích hoạt; tóm tắt hiện ra tức thì trên giao diện. \\
\hline
\textbf{Luồng phát sinh} & Gemini API không phản hồi $\to$ job thất bại, trạng thái chuyển về \textit{failed}; tác nhân có thể thử lại. \\
\hline
\textbf{Hậu điều kiện} & Văn bản tóm tắt xuất hiện trên ProjectDetailPage; ban quản lý nắm bắt nhanh kết quả chuyến đi. \\
\hline
\end{longtable}

\end{document}
